\subsection{Operator Classification}

\begin{perspective} \label{per:operator-module-similarity-dictionary} \hlt{Operator Modules and Similarity}\\
For $T\in\LL(\mathcal{V})$, $\mathcal{V}_T$ has action $p\cdot\mathbf{v}=p(T)\mathbf{v}$. Then $S\sim T\Longleftrightarrow\mathcal{V}_S\cong_{\mathbb{F}[z]}\mathcal{V}_T$; $\mathbb{F}[z]/\langle f\rangle\Longleftrightarrow C_f$ (Theorem \ref{thm:cyclic-subspace-quotient-and-companion}).
\end{perspective}

\begin{theorem} \label{thm:elementary-divisors-complete-similarity-invariant} \hlt{Complete Similarity Invariants}\\
Let $T,S\in\LL(\mathcal{V})$ and $A,B\in\operatorname{Mat}(n,n,\mathbb{F})$.
\begin{enumerate}[label=(\roman*)]
\setlength{\itemsep}{0pt}
\item $T\sim S\Longleftrightarrow\mathcal{V}_T\cong_{\mathbb{F}[z]}\mathcal{V}_S\Longleftrightarrow\operatorname{ElemDiv}(T)=\operatorname{ElemDiv}(S)\Longleftrightarrow\operatorname{InvFact}(T)=\operatorname{InvFact}(S)$.
\item $A\sim B\Longleftrightarrow(\mathbb{F}^n)_A\cong_{\mathbb{F}[z]}(\mathbb{F}^n)_B\Longleftrightarrow\operatorname{ElemDiv}(A)=\operatorname{ElemDiv}(B)\Longleftrightarrow\operatorname{InvFact}(A)=\operatorname{InvFact}(B)$.
\item If $A=[T]_{\mathcal{B}}$ for some basis $\mathcal{B}$ of $\mathcal{V}$, then $\operatorname{ElemDiv}(A)=\operatorname{ElemDiv}(T)$ and $\operatorname{InvFact}(A)=\operatorname{InvFact}(T)$.
\end{enumerate}
\begin{proof}
\begin{enumerate}[label=(\roman*).]
\setlength{\itemsep}{0pt}
\item[$(i)$]($\Rightarrow$) If $T\sim S$, then $\mathcal{V}_T\cong_{\mathbb{F}[z]}\mathcal{V}_S$ by Theorem \ref{thm:operator-similarity-iff-module-isomorphism}.\\
Module isomorphism preserves the PID decomposition, hence preserves $\operatorname{ElemDiv}$ and $\operatorname{InvFact}$.
\item[$(i)$]($\Leftarrow$) Equality of either $\operatorname{ElemDiv}$ or $\operatorname{InvFact}$ gives $\mathcal{V}_T\cong_{\mathbb{F}[z]}\mathcal{V}_S$ by uniqueness of the PID decomposition.\\
Thus $T\sim S$ by Theorem \ref{thm:operator-similarity-iff-module-isomorphism}.
\item[$(ii)$]($\Rightarrow$) If $A\sim B$, then $(\mathbb{F}^n)_A\cong_{\mathbb{F}[z]}(\mathbb{F}^n)_B$ by Corollary \ref{cor:matrix-similarity-iff-matrix-modules-isomorphic}; hence the divisor data agree.
\item[$(ii)$]($\Leftarrow$) Equality of either matrix divisor datum gives $(\mathbb{F}^n)_A\cong_{\mathbb{F}[z]}(\mathbb{F}^n)_B$, so $A\sim B$ by Corollary \ref{cor:matrix-similarity-iff-matrix-modules-isomorphic}.
\item[$(iii)$] If $A=[T]_{\mathcal{B}}$, then $\Phi(\mathbf{v})=[\mathbf{v}]_{\mathcal{B}}$ satisfies $\Phi(p(T)\mathbf{v})=p(A)\Phi(\mathbf{v})$.\\
Hence $\Phi:\mathcal{V}_T\to(\mathbb{F}^n)_A$ is an $\mathbb{F}[z]$-isomorphism, so the divisor data agree.
\end{enumerate}
\hlt{Structural Commentary}: Similarity is module isomorphism plus PID data.\\
\hlt{Key Insight}: Elementary divisors classify operator modules.
\end{proof}
\end{theorem}

\begin{lemma} \label{lem:determinant-of-companion-matrix} \hlt{Determinant of a Companion Block}\\
Let $f(z)=z^n+\sum_{i=0}^{n-1}c_iz^i$ be monic. Then $\det(zI_n-C_f)=f(z)$.
\begin{proof} Set $D_n(c_0,\ldots,c_{n-1})\defeq\det(zI_n-C_f)$. For $n=1$, $D_1(c_0)=z+c_0$.\\
For $n\geq2$, first-row expansion gives $D_n(c_0,\ldots,c_{n-1})=zD_{n-1}(c_1,\ldots,c_{n-1})+c_0$, since the $(1,1)$-minor is again companion-shaped and the last first-row cofactor contributes $c_0$. Induction gives $D_n=z^n+\sum_{i=0}^{n-1}c_iz^i=f(z)$.\\
\hlt{Structural Commentary}: The last column of $C_f$ stores the coefficients of the polynomial relation.\\
\hlt{Key Insight}: The shifted determinant reads that stored relation back as $f$.
\end{proof}
\end{lemma}

\begin{definition} \label{def:rational-canonical-block-forms} \hlt{Rational Canonical Block Forms}
\begin{enumerate}[label=(\roman*)]
\setlength{\itemsep}{0pt}
\item \hlt{Invariant-factor RCF}: $A=\operatorname{diag}(C_{s_1},\ldots,C_{s_r})$, where $s_j\in\mathbb{F}[z]$ are monic nonconstant and $s_r\mid\cdots\mid s_1$.
\item \hlt{Elementary-divisor RCF}: $A=\operatorname{diag}(C_{q_1},\ldots,C_{q_m})$, where each $q_j$ is power of monic irreducible polynomial.
\end{enumerate}
\end{definition}

\begin{theorem} \label{thm:rational-canonical-form-skeleton} \hlt{Rational Canonical Form Skeleton}\\
For $A\in\operatorname{Mat}(n,n,\mathbb{F})$, there is a unique tuple $\operatorname{InvFact}(A)=(s_1,\ldots,s_r)$ with $s_j\in\mathbb{F}[z]$ monic nonconstant, $s_r\mid\cdots\mid s_1$, and $A\sim_{\mathbb{F}}R_A\defeq\operatorname{diag}(C_{s_1},\ldots,C_{s_r})$. Moreover, $\mu_A=s_1$ and $p_A=\prod_{j=1}^r s_j$.\\
Equivalently, for $A,B\in\operatorname{Mat}(n,n,\mathbb{F})$, $A\sim_{\mathbb{F}}B\Longleftrightarrow\operatorname{ElemDiv}(A)=\operatorname{ElemDiv}(B)$ as multisets.
\begin{proof} Let $T_A:\mathbb{F}^n\to\mathbb{F}^n$, $\mathbf{x}\mapsto A\mathbf{x}$.\\
Theorem \ref{thm:operator-invariant-factor-skeleton} gives $\mathbb{F}^n=\bigoplus_{j=1}^r\mathcal{C}_j$, $\mathcal{C}_j\cong_{\mathbb{F}[z]}\mathbb{F}[z]/\langle s_j\rangle$, $s_r\mid\cdots\mid s_1$, $[T_A]=R_A$, $\mu_{T_A}=s_1$, and $p_{T_A}=\prod_{j=1}^r s_j$.\\
Hence $A\sim_{\mathbb{F}}R_A$, $\mu_A=s_1$, and $p_A=\prod_{j=1}^r s_j$; uniqueness of $\operatorname{InvFact}(A)$ is uniqueness in the same theorem.
\begin{enumerate}[label=\roman*.]
\setlength{\itemsep}{0pt}
\item[$(\Rightarrow)$]
If $A\sim_{\mathbb{F}}B$, then Theorem \ref{thm:elementary-divisors-complete-similarity-invariant} gives $\operatorname{ElemDiv}(A)=\operatorname{ElemDiv}(B)$.
\item[$(\Leftarrow)$]
If $\operatorname{ElemDiv}(A)=\operatorname{ElemDiv}(B)$, then Theorem \ref{thm:elementary-divisors-complete-similarity-invariant} gives $A\sim_{\mathbb{F}}B$.
\end{enumerate}
\hlt{Structural Commentary}: Companion blocks classify similarity without requiring $p_A$ to split.\\
\hlt{Key Insight}: Similarity classes can be represented by polynomial data rather than by choosing roots.
\end{proof}
\end{theorem}

\begin{theorem} \label{thm:rational-canonical-form-elementary-and-invariant-versions} \hlt{Rational Canonical Form: Two Versions}\\
Let $T\in\LL(\mathcal{V})$.
\begin{enumerate}[label=(\roman*)]
\setlength{\itemsep}{0pt}
\item If $\operatorname{ElemDiv}(T)=\{q_1,\ldots,q_m\}$, then $\exists$ basis $\mathcal{R}$ with $[T]_{\mathcal{R}}=\operatorname{diag}(C_{q_1},\ldots,C_{q_m})$.\\
The form is unique up to block permutation.
\item If $\operatorname{InvFact}(T)=(s_1,\ldots,s_r)$ and $s_r\mid\cdots\mid s_1$, then $\exists$ basis $\mathcal{B}$ with $[T]_{\mathcal{B}}=\operatorname{diag}(C_{s_1},\ldots,C_{s_r})$.\\
The form is unique in this order.
\item If $A\in\operatorname{Mat}(n,n,\mathbb{F})$, then the similarity class $[A]_{\sim_{\mathbb{F}}}$ contains both forms.\\
Elementary-divisor form is unique up to block order, invariant-factor form by its divisibility chain.
\end{enumerate}
\begin{proof}
\begin{enumerate}[label=(\roman*)]
\setlength{\itemsep}{0pt}
\item Theorem \ref{thm:operator-primary-cyclic-decomposition} gives $\mathcal{V}=\bigoplus_{\alpha=1}^m\mathcal{C}_\alpha$ with orders $q_\alpha$; Theorem \ref{thm:cyclic-subspace-quotient-and-companion} gives cyclic bases and blocks $C_{q_\alpha}$.\\
Permuting summands conjugates by a permutation matrix. 
\item Theorem \ref{thm:operator-invariant-factor-skeleton} gives $\mathcal{V}=\bigoplus_{j=1}^r\mathcal{C}_j$, $\mathcal{C}_j\cong_{\mathbb{F}[z]}\mathbb{F}[z]/\langle s_j\rangle$, and cyclic bases with block matrix $\operatorname{diag}(C_{s_1},\ldots,C_{s_r})$.
\item For $A$, apply (i),(ii) to $T_A:\mathbf{x}\mapsto A\mathbf{x}$; any chosen basis gives a similarity matrix.\\
Uniqueness follows from Theorem \ref{thm:elementary-divisors-complete-similarity-invariant}.
\end{enumerate}
\hlt{Structural Commentary}: Elementary divisors split primary blocks; invariant factors recombine them.\\
\hlt{Key Insight}: Both forms classify the same similarity classes, but they expose different polynomial data.
\end{proof}
\end{theorem}

\begin{method} \label{meth:construct-operator-from-elementary-divisors} \hlt{Constructing an Operator from Elementary Divisors}\\
Let $Q=\{q_1,\ldots,q_m\}$ be a finite multiset of monic prime powers in $\mathbb{F}[z]$, and set $n\defeq\sum_{j=1}^m\deg q_j$.
\begin{enumerate}[label=(\roman*)]
\setlength{\itemsep}{0pt}
\item Set $R_Q\defeq\operatorname{diag}(C_{q_1},\ldots,C_{q_m})\in\operatorname{Mat}(n,n,\mathbb{F})$.
\item Then $\operatorname{ElemDiv}(T_{R_Q})=Q$; hence every such $Q$ is realised on $\mathbb{F}^n$.
\item Apply Method \ref{meth:elementary-divisors-to-invariant-factors} to recombine $Q$ into $\operatorname{InvFact}(T_{R_Q})=(s_1,\ldots,s_r)$, then $\operatorname{diag}(C_{s_1},\ldots,C_{s_r})\sim_{\mathbb{F}}R_Q$.
\end{enumerate}
\hlt{Structural Commentary}: RCF constructs operators from divisor data.\\
\hlt{Key Insight}: Companion blocks realise cyclic summands; block sums realise direct sums.
\end{method}

\begin{example} \label{ex:rational-canonical-forms-from-minimal-polynomial} \hlt{Rational Canonical Possibilities from Degree Counts}\\
Let $T\in\LL(\mathbb{R}^7)$ and $\mu_T=(z-1)(z^2+1)^2$. For $Q=\operatorname{ElemDiv}(T)$, $\operatorname{lcm}Q=\mu_T$ and $\sum_{q\in Q}\deg q=7$.\\
Hence $Q_1=\{z-1,(z^2+1)^2,z^2+1\}$ or $Q_2=\{z-1,z-1,z-1,(z^2+1)^2\}$.\\
Thus the elementary-divisor RCF is $R_i=\bigoplus_{q\in Q_i}C_q$ for $i=1,2$, up to block order.
\end{example}

\begin{theorem} \label{thm:determinant-form-characteristic-polynomial} \hlt{Determinant Form of the Characteristic Polynomial}\\
For $A\in\operatorname{Mat}(n,n,\mathbb{F})$, $\chi_A(z)=\det(zI_n-A)$. If $A=[T]_{\mathcal{B}}$, then $\chi_T=\chi_A$.
\begin{proof}
Let $Q\defeq\operatorname{ElemDiv}(A)$ and choose elementary-divisor RCF $R=\bigoplus_{q\in Q}C_q\sim_{\mathbb{F}}A$.\\
By Lemma \ref{lem:determinant-of-companion-matrix}, $\det(zI-R)=\prod_{q\in Q}\det(zI-C_q)=\prod_{q\in Q}q=\chi_R$.\\
Since similarity preserves both $Q$ and $\det(zI-\cdot)$, $\chi_A=\chi_R=\det(zI-R)=\det(zI-A)$.\\
If $A=[T]_{\mathcal{B}}$, then $\operatorname{ElemDiv}(A)=\operatorname{ElemDiv}(T)$, hence $\chi_A=\chi_T$.\\
\hlt{Structural Commentary}: Rational canonical form reduces the determinant computation to companion blocks.\\
\hlt{Key Insight}: The abstract product of elementary divisors is the concrete shifted determinant.
\end{proof}
\end{theorem}

\begin{theorem} \label{thm:base-field-change-elementary-divisors} \hlt{Changing the Base Field}\\
Let $\mathbb{F}\subseteq\mathbb{K}$, $A\in\operatorname{Mat}(n,n,\mathbb{F})$, and $\operatorname{ElemDiv}_{\mathbb{F}}(A)=\{p_i^{e_{i,j}}\}_{i,j}$ with $p_i\in\mathbb{F}[z]$ monic irreducible.\\
If $p_i=\prod_{\ell=1}^{m_i}a_{i,\ell}^{d_{i,\ell}}$ in $\mathbb{K}[z]$, with $a_{i,\ell}$ pairwise distinct monic irreducible over $\mathbb{K}$, then $\operatorname{ElemDiv}_{\mathbb{K}}(A)=\{a_{i,\ell}^{d_{i,\ell}e_{i,j}}\}_{i,j,\ell}$.\\
Moreover, $\operatorname{InvFact}_{\mathbb{K}}(A)=\operatorname{InvFact}_{\mathbb{F}}(A)$.
\begin{proof}
Take elementary-divisor RCF over $\mathbb{F}$: $R=\bigoplus_{i,j}C_{p_i^{e_{i,j}}}\sim_{\mathbb{F}}A$.\\
Since $R\in\operatorname{Mat}(n,n,\mathbb{K})$ and $p_i^{e_{i,j}}=\prod_{\ell=1}^{m_i}a_{i,\ell}^{d_{i,\ell}e_{i,j}}$ with coprime factors in $\mathbb{K}[z]$,\\
Theorem \ref{thm:cyclic-composition-coprime-orders} gives $C_{p_i^{e_{i,j}}}\sim_{\mathbb{K}}\bigoplus_{\ell=1}^{m_i}C_{a_{i,\ell}^{d_{i,\ell}e_{i,j}}}$. Summing over $(i,j)$ yields $\operatorname{ElemDiv}_{\mathbb{K}}(A)$.\\
Recombination into the divisibility chain gives the original invariant factors.\\
\hlt{Structural Commentary}: Field extension refines elementary divisors.\\
\hlt{Key Insight}: Elementary divisors split; invariant factors persist.
\end{proof}
\end{theorem}

\begin{corollary} \label{cor:nonderogatory-iff-companion-similar} \hlt{Nonderogatory Means One Companion Block}\\
For $A\in\operatorname{Mat}(n,n,\mathbb{F})$, the following are equivalent:
\begin{enumerate}[label=(\roman*)]
\setlength{\itemsep}{0pt}
\item $A$ is nonderogatory, i.e. $\mu_A=\chi_A$;
\item $(\mathbb{F}^n)_A$ is cyclic;
\item $A\sim_{\mathbb{F}}C_{\chi_A}$.
\end{enumerate}
\begin{proof}
\begin{enumerate}[label=(\roman*).]
\setlength{\itemsep}{0pt}
\item[$(i)$]$\Rightarrow (ii)$ This is Theorem \ref{thm:nonderogatory-cyclic-equivalences}.
\item[$(ii)$]$\Rightarrow (iii)$ Cyclicity gives $\operatorname{InvFact}(A)=(s)$. By Theorem \ref{thm:operator-invariant-factor-skeleton}, $s=\mu_A=\chi_A$; hence $A\sim_{\mathbb{F}}C_{\chi_A}$.
\item[$(iii)$]$\Rightarrow (i)$ Since $\mu_{C_f}=f$ by Theorem \ref{thm:companion-matrix-has-prescribed-minimal-polynomial}, similarity gives $\mu_A=\mu_{C_{\chi_A}}=\chi_A$.
\end{enumerate}
\hlt{Structural Commentary}: Nonderogatory means one cyclic block.\\
\hlt{Key Insight}: One cyclic generator gives one companion block.
\end{proof}
\end{corollary}

\begin{corollary} \label{cor:diagonal-rational-canonical-form-criterion} \hlt{When the Elementary-Divisor Form Is Diagonal}\\
Let $T\in\LL(\mathcal{V})$, $Q=\operatorname{ElemDiv}(T)$, and $R_{\operatorname{ed}}=\bigoplus_{q\in Q}C_q$.\\
Then $R_{\operatorname{ed}}$ is diagonal iff $\deg q=1$ for all $q\in Q$, iff $\mu_T=\prod_{\lambda\in\Lambda}(z-\lambda)$ for some finite $\Lambda\subseteq\mathbb{F}$.
\begin{proof}
\begin{enumerate}[label=\roman*.]
\setlength{\itemsep}{0pt}
\item[$(\Rightarrow)$] If $R_{\operatorname{ed}}$ is diagonal, each block $C_q$ has size $1$, hence $\deg q=1$ for all $q\in Q$.\\
Thus $\mu_T=\operatorname{lcm}Q$ is a product of distinct linear factors.
\item[$(\Leftarrow)$] If $\mu_T=\prod_{\lambda\in\Lambda}(z-\lambda)$, then each prime power $q\in Q$ divides $\mu_T$, so $q=z-\lambda$ for some $\lambda\in\Lambda$.\\
Hence each $C_q$ is $1\times1$, and $R_{\operatorname{ed}}$ is diagonal.
\end{enumerate}
\hlt{Structural Commentary}: Diagonal RCF means all companion blocks are $1\times1$.\\
\hlt{Key Insight}: Nonlinear or repeated factors obstruct diagonal RCF.
\end{proof}
\end{corollary}

\begin{proposition} \label{prop:ab-ba-characteristic-polynomial-and-similarity} \hlt{$AB$ and $BA$ Have the Same Characteristic Polynomial}\\
For $A,B\in\operatorname{Mat}(n,n,\mathbb{F})$, $\chi_{AB}=\chi_{BA}$.\\
If $A\in\operatorname{GL}(n,\mathbb{F})$ or $B\in\operatorname{GL}(n,\mathbb{F})$, then $AB\sim_{\mathbb{F}}BA$; otherwise not necessarily.
\begin{proof}
In $\mathbb{F}(z)$, by Sylvester's identitym $\det(zI-AB)=z^n\det(I-z^{-1}AB)=z^n\det(I-z^{-1}BA)=\det(zI-BA)$, hence $\chi_{AB}=\chi_{BA}$. If $A$ is invertible, $BA=A^{-1}(AB)A$; if $B$ is invertible, $AB=B^{-1}(BA)B$.\\
Without invertibility, take $A=\begin{bmatrix}1&0\\0&0\end{bmatrix}$ and $B=\begin{bmatrix}0&1\\0&0\end{bmatrix}$.\\
Then $AB=B$, $BA=0$, $\chi_{AB}=\chi_{BA}=z^2$, but $\operatorname{rank}(AB)=1\neq0=\operatorname{rank}(BA)$, so $AB\not\sim_{\mathbb{F}}BA$.\\
\hlt{Structural Commentary}: Characteristic polynomials see less than similarity.\\
\hlt{Key Insight}: $AB$ and $BA$ share determinant data; invertibility yields similarity.
\end{proof}
\end{proposition}

\begin{example} \label{ex:base-field-rational-canonical-form-minimal-polynomial} \hlt{The Same Polynomial over Different Fields}\\
Let $\dim V=4$ and $\mu_T=(z^2+1)(z^2-2)$. Since $\deg\mu_T=4$, $T$ is nonderogatory and $\operatorname{InvFact}_K(T)=(\mu_T)$ for $K=\mathbb{Q},\mathbb{R},\mathbb{C}$; hence $R_{\operatorname{inv},K}=C_{\mu_T}$. Meanwhile
\begin{center}
\begin{tabular}{c|c}
$K$ & $\operatorname{ElemDiv}_K(T)$\\
\hline
$\mathbb{Q}$ & $\{z^2+1,\ z^2-2\}$\\
$\mathbb{R}$ & $\{z^2+1,\ z-\sqrt{2},\ z+\sqrt{2}\}$\\
$\mathbb{C}$ & $\{z-i,\ z+i,\ z-\sqrt{2},\ z+\sqrt{2}\}$
\end{tabular}
\end{center}
Thus $R_{\operatorname{ed},K}=\bigoplus_{q\in\operatorname{ElemDiv}_K(T)}C_q$: scalar extension refines elementary divisors while preserving the invariant factor. Over $\mathbb{C}$, $R_{\operatorname{ed},\mathbb{C}}$ is diagonal.
\end{example}

\begin{figure}[H]
  \centering
  \begin{tikzpicture}[
  >=Stealth,
  every node/.style={font=\small},
  class/.style={draw=blue!60!black, fill=blue!7, rounded corners=2pt,
    minimum width=2.25cm, minimum height=0.72cm, align=center, inner sep=5pt},
  data/.style={draw=dartmouthgreen!65!black, fill=dartmouthgreen!8, rounded corners=2pt,
    minimum width=2.9cm, minimum height=0.72cm, align=center, inner sep=5pt},
  form/.style={draw=gray!60, fill=gray!6, rounded corners=2pt,
    minimum width=3.45cm, minimum height=0.72cm, align=center, inner sep=5pt},
  link/.style={->, line width=0.75pt, draw=black!70},
  note/.style={font=\scriptsize, inner sep=2pt, fill=white}
]
  \node[class] (ops) at (0,0) {$[T]_{\sim}$};
  \node[class] (mods) at (4.3,0) {$[\mathcal{V}_T]_{\mathbb{F}[z]}$};
  \node[data] (invs) at (8.7,0) {$\{s_i\}:s_r\mid\cdots\mid s_1$\\$\{q_j\}$};
  \node[form] (mats) at (13.4,0) {$\operatorname{diag}(C_{s_1},\ldots,C_{s_r})$};

  \draw[link] (ops) -- node[note, above] {$p\cdot\mathbf{v}=p(T)\mathbf{v}$} (mods);
  \draw[link] (mods) -- node[note, above] {$\bigoplus_i\mathbb{F}[z]/\langle s_i\rangle$} (invs);
  \draw[link] (invs) -- node[note, above] {$s_i\mapsto C_{s_i}$} (mats);
\end{tikzpicture}

  \caption{Rational-canonical dictionary: $[T]_{\sim}$, $[\mathcal{V}_T]_{\mathbb{F}[z]}$, $\{s_i\}$ or $\{q_j\}$, and $\operatorname{diag}(C_{s_i})$ are equivalent data.}
  \label{fig:operator-module-canonical-form-dictionary}
\end{figure}
